%\ifodd\value{page}\hbox{}\newpage\fi
\chapter*{Ṛgveda 10.125 — Devī Sūkta, Mantra 4}
\addcontentsline{toc}{chapter}{Devī Sūkta 10.125.4}

\begin{sloka}
{\sanskritfont
\begin{tabular}{c}
अहं एव स्वयमिदं वदामि जुष्टं देवेभिरुत मानुषेभिः ।\\
यं कामये तं तमुग्रं कृणोमि तं ब्रह्माणं तमृषिं तं सुमेधाम् ॥४॥
\end{tabular}

\vspace{1.5em}

{\middlesize \iastfont
ahaṁ eva svayam idaṁ vadāmi juṣṭaṁ devebhir uta mānuṣebhiḥ |\\
yaṁ kāmaye taṁ tam ugraṁ kṛṇomi taṁ brahmāṇaṁ tam ṛṣiṁ taṁ sumedhām ||4||}

\vspace{1em}

{\middlesize
\anvaya{%
अहं एव स्वयम् इदम् वदामि ।
इदम् देवेभिः उत मानुषेभिः जुष्टम् ।
यम् कामये तम् तम् उग्रम् कृणोमि ।
तम् ब्रह्माणम् तम् ऋषिम् तम् सुमेधाम् ।
}
}

\middlesize \iastfont \itmean{%
«Io sola, da me stessa, proclamo questo,  
ciò che è caro agli Dei e agli uomini.  
Colui che io desidero, lo rendo potente;  
lo rendo sacerdote, veggente, dotato di profonda intelligenza.»%
}

}
\end{sloka}

\wordmeaning{%
\wordline{अहं एव स्वयम् इदम् वदामि}
{ahaṁ eva svayam idaṁ vadāmi}
{ io stessa (\textit{aham eva svayam}) proclamo (\textit{vadāmi}) questo; }%

\wordline{देवेभिः उत मानुषेभिः जुष्टम्}
{devebhiḥ uta mānuṣebhiḥ juṣṭam}
{ ciò che è gradito, caro (\textit{juṣṭa}) agli Dei e anche agli uomini; }%

\wordline{यम् कामये तम् उग्रम् कृणोमि}
{yaṁ kāmaye taṁ ugram kṛṇomi}
{ colui che io desidero (\textit{kāmaye}), lo rendo (\textit{kṛṇomi}) potente, energico (\textit{ugra}); }%

\wordline{तम् ब्रह्माणम् तम् ऋषिम् तम् सुमेधाम्}
{taṁ brahmāṇaṁ taṁ ṛṣiṁ taṁ sumedhām}
{ lo rendo sacerdote (\textit{brahmāṇa}), veggente (\textit{ṛṣi}) e dotato di eccellente intelligenza (\textit{sumedhā}); }%
}

Nel quarto mantra la Dea afferma in modo esplicito la propria sovranità sulla parola e sull’efficacia. Ella non trasmette un messaggio ricevuto: parla “da sé stessa”. Ciò che proclama è detto essere gradito tanto agli Dei quanto agli uomini, indicando una verità che non appartiene a un solo ordine, ma li attraversa entrambi. La seconda metà del mantra introduce una dimensione decisiva: la Dea non solo parla, ma conferisce potenza. Colui che ella desidera viene reso \textit{ugra}, cioè capace di forza ed efficacia; diventa \textit{brahmāṇa}, \textit{ṛṣi} e \textit{sumedhā}. Autorità rituale, visione ispirata e intelligenza penetrante non sono qualità innate o acquisite autonomamente: sono partecipazioni alla sua energia. Questo mantra chiarisce così che la conoscenza, la parola autorevole e la capacità di guidare non derivano semplicemente dallo sforzo umano. Esse scaturiscono da una relazione con la Dea, che sceglie, investe e rende capace. La sovranità divina si manifesta qui come potere di far emergere la competenza spirituale e l’intelligenza creativa.

\textit{Svayam} letteralmente significa “da sé”, “per se stessa”, ma il suo valore qui è eminentemente teologico e non semplicemente riflessivo. Quando la Dea afferma \textit{ahaṁ eva svayam idaṁ vadāmi}, dichiara che la sua parola non dipende da alcuna fonte esterna. Non è trasmissione, non è eco, non è rivelazione mediata; è parola che nasce da sé stessa e in sé stessa trova il proprio fondamento. \textit{Svayam} indica dunque una autofondazione della voce sacra, una parola che non chiede legittimazione perché è essa stessa principio di legittimità. In questo senso, la sovranità della Dea non è delegata né derivata: è originaria. Da qui discende immediatamente la sua efficacia (\textit{kṛṇomi}): proprio perché la parola nasce da sé, essa può trasformare, rendere \textit{ugra}, generare \textit{brahmāṇa}, \textit{ṛṣi} e \textit{sumedhā}. Il mantra mostra così che autorità, conoscenza e potere creativo non hanno la loro radice nello sforzo umano o nella tradizione ricevuta, ma in una sorgente che si dà da sé. \textit{Svayam} nomina quindi il cuore della sovranità divina: una presenza che parla, agisce e istituisce senza bisogno di altro che di sé stessa.
